\section{An elementary class proof for the conductor-seven residual field}
\label{sec:real7-cm-elementary-class}

\paragraph{Scope.}
This appendix supplies an elementary proof of the fixed class input
in (G184.3).  It is separate from the frozen quotient-certificate
proof used there.  It adds no character family, weight, coefficient
prime, or common-window conclusion.

\begin{theorem}[AC-G185: the fixed degree-six field has class number one]
\label{AC-G185}
Let $\omega^2+\omega+1=0$, let $\zeta$ be a primitive seventh
root of unity, and put
\[
 K=\mathbb Q(\omega),\qquad
 a=\zeta+\zeta^{-1},\qquad
 F=\mathbb Q(a),\qquad L=KF.
 \tag{G185.1}
\]
Then $\operatorname{Cl}(L)=0$.  The proof uses the discriminant
of an explicit integral order as an upper bound and does not require
a computed maximal order, class group, regulator, or unit basis.
In particular it independently supplies the class-group input of
AC-G184, whose original coefficient and local-condition statements
are unchanged.
\end{theorem}

\begin{proof}
\proofstep{1. Degree and an integral-order discriminant.}
The relation $1+\zeta+\cdots+\zeta^6=0$ gives
\[
 f(a)=0,\qquad f(Y)=Y^3+Y^2-2Y-1.
 \tag{G185.2}
\]
Neither $1$ nor $-1$ is a root of this monic cubic, so it is
irreducible over $\mathbb Q$.  Its three roots are real, namely
$\zeta^j+\zeta^{-j}$ for $j=1,2,3$.
Thus $[F:\mathbb Q]=3$, $K\cap F=\mathbb Q$, and $L$ has
degree six and signature $(0,3)$.

All six elements $1,a,a^2,\omega,a\omega,a^2\omega$ are integral
and form a rational basis of $L$.  Their integral span is the order
$A=\mathbb Z[a,\omega]$.  In these product bases the trace pairing
is the tensor product of the two trace pairings.  The latter have
matrices
\[
 \begin{pmatrix}2&-1\\-1&-1\end{pmatrix},\qquad
 \begin{pmatrix}3&-1&5\\-1&5&-4\\5&-4&13\end{pmatrix}
 \tag{G185.3}
\]
and determinants $-3$ and $49$.  Therefore
\[
 \operatorname{disc}(A)=(-3)^3\,49^2=-64827,
 \qquad
 |D_L|=\frac{64827}{[\mathcal O_L:A]^2}\leq64827.
 \tag{G185.4}
\]
Only the inequality is used; equality with the field discriminant
is not needed.

\proofstep{2. An exact bound below ten.}
Minkowski's ideal-class bound gives an integral ideal in each class
whose norm is at most
\[
 B_L=\left(\frac4\pi\right)^3\frac{6!}{6^6}\sqrt{|D_L|}.
 \tag{G185.5}
\]
Here the convention is the ordinary ideal class group, and the
three complex places contribute the exponent three.  Since
$\pi>3$, $64827<255^2$, and $6!/6^6=5/324$, we have the exact
rational estimate
\[
 B_L<\frac{64}{27}\,\frac5{324}\,255
       =\frac{20400}{2187}<10.
 \tag{G185.6}
\]
Every ideal class consequently has an integral representative of
integer norm at most nine.

\proofstep{3. Every prime ideal of norm at most nine lies above seven.}
For completeness, the values of $f$ at all elements of the three
small residue fields are
\[
\begin{array}{c|c|c}
\ell&\text{arguments}&\text{values of }f\pmod\ell\\ \hline
2&0,1&1,1\\
3&0,1,2&2,2,1\\
5&0,1,2,3,4&4,4,2,4,1.
\end{array}
 \tag{G185.7}
\]
Thus $f$ is irreducible modulo $2,3,5$.  Its discriminant $49$
is prime to these primes, so each gives the unramified cubic
extension in $F$.  The polynomial $Y^2+Y+1$ is irreducible modulo
$2$ and $5$, so $K$ gives their unramified quadratic extensions.
Their local composita have unramified degree six.  Hence the norms
of primes of $L$ above $2$ and $5$ are $2^6$ and $5^6$.

At $3$, the field $K=\mathbb Q(\sqrt{-3})$ is ramified quadratic:
$Y^2+3$ is Eisenstein over $\mathbb Q_3$.  The unramified cubic
extension supplied by $F$ remains linearly disjoint from this
quadratic extension.  Its compositum therefore has ramification
index two and residue degree three.  Every prime of $L$ above
$3$ has norm $3^3=27$.  In particular this ramified rational prime
cannot contribute an ideal of norm at most nine.

At $7$, the translation
\[
                   f(Y+2)=Y^3+7Y^2+14Y+7
 \tag{G185.8}
\]
is Eisenstein, so $F$ is totally ramified of degree three.
The two distinct roots $2,4$ of $Y^2+Y+1$ modulo seven lift to
$\mathbb Q_7$, so $K$ splits there.  There are exactly two primes
of $L$ above seven, each of residue degree one and norm seven.
These are the only possible prime ideals of norm at most nine,
since the rational primes at most nine are $2,3,5,7$.

\proofstep{4. Explicit generators for the two norm-seven primes.}
Taking the product over the three conjugates of $a$ gives
\[
 N_{L/K}(a-\omega)=-f(\omega)=1+3\omega,
 \qquad N_{K/\mathbb Q}(1+3\omega)=7.
 \tag{G185.9}
\]
Indeed $f(\omega)=\omega^2-2\omega=-1-3\omega$, and
$N_{K/\mathbb Q}(x+y\omega)=x^2-xy+y^2$.
Thus the integral principal ideal $(a-\omega)$ has norm seven
and is a prime ideal.  Complex conjugation fixes $a$ and sends
$\omega$ to $\omega^2$, so $(a-\omega^2)$ is another principal
prime of norm seven.

They are distinct.  Their ideal norms to $K$ are respectively
$(1+3\omega)$ and $(1+3\omega^2)$, and
\[
                  (1+3\omega)+(1+3\omega^2)=-1.
 \tag{G185.10}
\]
Consequently these two norm ideals are coprime.  Each has norm
seven, and the relative residue degrees are one, so they are the
distinct contractions of the two primes above seven.  The two
principal ideals just constructed therefore exhaust those primes.

An integral ideal of norm at most nine either is the unit ideal
or has one of these prime ideals as a factor.  A product of two
nontrivial such factors already has norm at least $7^2>9$.
It follows that every representative from Step~2 is principal,
which proves $\operatorname{Cl}(L)=0$.
\end{proof}

\paragraph{Source and dependency scope.}
The sole general ideal-class estimate is Milne,
\emph{Algebraic Number Theory}, v3.08 (19 July 2020),
Theorem~4.3, printed p.~70.  Its complete proof on printed
p.~81 and the applied lattice, volume and short-element proofs
(Theorems~4.17 and~4.19, Lemmas~4.22--4.23,
Propositions~4.26--4.27, printed pp.~75--81) were freshly read.
This is the same Minkowski input already used in Step~1 of
the finite-gluing estimate.  Basic ideal factorization, norm
multiplicativity, Hensel lifting and unramified local extensions
retain their usual algebraic-number-theory scope.  All field,
discriminant, small-prime and generator calculations needed here
are displayed above.  The appendix strengthens a fixed dependency
of AC-G184 and makes no further reduction of the common-window
or all-prime main gate.
